\section{Introduction}
\label{sec:introduction}

Fix \(\nu>0\) and a divergence-free datum
\(u_0\in H^\infty(\R^3;\R^3)\).  Smooth-data local theory gives a unique
classical solution on a maximal interval \([0,T_*)\)
\cite{FujitaKato1964,Kato1984}.  Throughout the paper, \((u,p)\) denotes that
solution.  The problem is to decide whether \(T_*\) can be finite.  On each strict
subinterval \([0,T]\), smoothness already makes every finite Sobolev norm
finite.  That fact alone says nothing about what happens as \(T\uparrow T_*\):
the bound may depend on \(T\) and diverge at the endpoint.  The proof must
instead control the approach to \(T_*\) with constants that do not depend on
the cutoff.

Leray constructed global finite-energy weak solutions \cite{Leray1934}, and
Caffarelli, Kohn, and Nirenberg bounded the size of the possible singular set
for suitable weak solutions \cite{CKN1982}.  Serrin proved regularity under
scale-balanced spacetime integrability hypotheses \cite{Serrin1962};
Escauriaza, Seregin, and \v{S}ver\'ak treated the endpoint
\(L^\infty_tL^3_x\) case \cite{ESS2003}.  Fujita--Kato and Kato supply the
strong local theory used below at the terminal restart
\cite{FujitaKato1964,Kato1984}, and Koch--Tataru place small-data
global well-posedness in the critical space \(BMO^{-1}\)
\cite{KochTataru2001}.
Tao gives a localization and compactness formulation of the global problem
\cite{Tao2013}.  The regularity criteria show that control of an appropriate
critical quantity excludes singularity.  The argument below does not assume
such control.  It derives it from one family of estimates for all finite time
and space derivatives of the maximal solution.

Pressure is part of that derivative structure.  With the decay normalization,
incompressibility determines it from the current velocity:
\[
 -\Delta p=\partial_i\partial_j(u_i u_j),
 \qquad
 p=R_iR_j(u_i u_j).
\]
Here \(R_i\) are the Riesz transforms, and repeated indices are summed.
Write
\[
 V_n:=\partial_t^nu,
 \qquad
 Q_n:=\partial_t^np.
\]
Differentiating the equation does not introduce a new solution.  It gives
\begin{align}
 Q_n
 &=R_iR_j\sum_{a=0}^{n}\binom naV_{a,i}V_{n-a,j},
 \label{eq:intro-pressure-recurrence}\\
 V_{n+1}
 &=\nu\Delta V_n
 -\sum_{a=0}^{n}\binom na(V_a\cdot\nabla)V_{n-a}
 -\nabla Q_n.
 \label{eq:intro-velocity-recurrence}
\end{align}
Thus the time derivatives, spatial derivatives, and pressure derivatives are
coordinates of one solution.  A time derivative raises the order of the
viscous term by two spatial derivatives and distributes the transport and
pressure terms among earlier time derivatives.  The resulting estimates must
track temporal depth and spatial depth independently.

For finite \(N,M\in\Nzero\), set
\begin{equation}
 \mathfrak R_{N,M}(u;T)
 :=\sum_{n=0}^{N}\norm{V_n}_{L^2(0,T;H^{M+1})}^2
 +\sum_{n=0}^{N}\norm{V_{n+1}}_{L^2(0,T;H^{M-1})}^2.
 \label{eq:intro-rectangle-norm}
\end{equation}
The central estimate is that, under the contradiction hypothesis
\(T_*<\infty\), every fixed finite pair satisfies
\begin{equation}
 \sup_{0<T<T_*}\mathfrak R_{N,M}(u;T)<\infty.
 \label{eq:intro-terminal-rectangle}
\end{equation}
This is Theorem~\ref{thm:datum-rectangles} below.
The bound is coordinatewise: its constant may depend on \(N\) and \(M\), but
it does not depend on \(T\).  No uniform estimate over all derivative orders
is asserted or needed.

Every rectangle contains derivatives of the same solution.  If
\(N'\le N\) and \(M'\le M\), restriction deletes the higher rows and applies
Sobolev embeddings to the rows that remain; a shared derivative does not
change.  Write \(\mathcal R_{N,M}[u_0;T_*]\) for the resulting whole-terminal
finite rectangle, and denote the compatible family by
\begin{equation}
 \mathcal I[u_0]
 :=\bigl(\mathcal R_{N,M}[u_0;T_*]\bigr)_{N,M\in\Nzero}.
 \label{eq:intro-compatible-intersection}
\end{equation}
Given finite \(r,m\), choose one rectangle deep enough to contain the
requested time and space derivatives.  Compatibility guarantees that this
coordinate is the same one in every larger rectangle, and hence
\begin{equation}
 u\in\bigcap_{r,m\in\Nzero}H^r(0,T_*;H^m).
 \label{eq:intro-all-orders-intersection}
\end{equation}

The zeroth time row of \eqref{eq:intro-all-orders-intersection} gives, for
each finite \(m\), \(u\in L^2(0,T_*;H^{m+2})\).  Sobolev multiplication gives
\[
 \norm{\Pp\nabla\!\cdot(u\otimes u)}_{H^m}
 \le C_m\norm u_{H^{m+2}}^2.
\]
Since \(T_*<\infty\), the nonlinear term and \(\Delta u\) belong to
\(L^1(0,T_*;H^m)\); for the linear term, explicitly,
\[
 \norm{\Delta u}_{L^1(0,T_*;H^m)}
 \le T_*^{1/2}\norm u_{L^2(0,T_*;H^{m+2})}.
\]
The equation
\[
 u_t=\nu\Delta u-\Pp\nabla\!\cdot(u\otimes u)
\]
now gives
\[
 u\in W^{1,1}(0,T_*;H^m)
 \hookrightarrow C([0,T_*];H^m)
 \qquad(m<\infty).
\]
The endpoint values agree under the Sobolev inclusions because they are limits
of the same function \(u(t)\).  They define one terminal state
\[
 u_*:=u(T_*)\in\bigcap_{m<\infty}H^m.
\]
Since \(u_*\in H^\infty_\sigma\), local well-posedness restarts the equation
at \(T_*\).  Equations \eqref{eq:intro-pressure-recurrence}--
\eqref{eq:intro-velocity-recurrence} recover the endpoint pressure and time
derivatives.  Uniqueness joins the restarted solution to \(u\), producing an
extension past \(T_*\).  This contradicts maximality and proves global
smoothness.

\clearpage
